\section{Scope, Method, and Qualifications}\label{sec:scope}

\subsection{Work identity}

\begin{center}
\small
\begin{tabular}{@{}>{\raggedright\arraybackslash}p{0.24\linewidth} >{\raggedright\arraybackslash}p{0.70\linewidth}@{}}
\toprule
Rite and books & Roman Rite; the \latin{Missale Romanum} of 1962 and the Roman Breviary as ordered by the Code of Rubrics of 1960 \\
Rubrics & \latin{Rubricae Breviarii et Missalis Romani}, \latin{Acta Apostolicae Sedis} 52 (1960), in force from 1 January 1961; cited by number \\
Season & Advent 2026: First Vespers of the First Sunday (Saturday 28 November) to None of the Vigil of the Nativity (Thursday 24 December) \\
Territory & The general Roman calendar. Law is stated for the Latin Church and, where the book says so, for the United States \\
Law as of & 29 November 2026, the First Sunday of Advent \\
Languages & English exposition; Latin from the editions named; English from the translations named \\
Reader & A Catholic lay reader who keeps or wishes to keep Advent with the 1962 books, and who reads for prayer \\
\bottomrule
\end{tabular}
\end{center}

\subsection{What is covered}

The meaning of the season as the Fathers and the medieval writers on the liturgy taught it; the
history of its observance, from the councils of Spain and Gaul and the December fast at Rome to
the books of 1962; the arrangement of the four weeks, the Ember Days, the last days and the Vigil
under the 1960 rubrics; the Sundays and Ember Days, each with its station, its governing texts at
the length of an opening phrase or a collect, its place in the season, and a reading from the
Fathers; the practice of the season in prayer, fasting, almsgiving, a rule of life and confession;
a reading for each day of Advent in any year, with 2026 as the worked year; and the prayers of the season, each from one edition.

\subsection{What is not covered}

\begin{itemize}[leftmargin=1.3em, itemsep=0.2em]
\item The complete Mass or Office of any day, which the Missal and Breviary give. For each Sunday
and proper day this book gives the station, a governing text or two by its opening words or, for a
collect, in full, the day's place in the season, and a reading from the Fathers; it prints no
complete proper and does not go through the parts of any Mass.
\item The Christmas season. The Vigil is treated as the last day of Advent; the feast of the
Nativity is not treated.
\item The feasts of the sanctoral cycle that fall in Advent, beyond their place in the daily table;
and any diocesan, national or religious calendar.
\item An Ordo. The daily table applies the general rubrics to the weeks and days of the season for
private reading and is not a statement of what any church will celebrate.
\item The particular law of any conference of bishops on holy days of obligation and Sunday Mass:
the book states the universal canons (cc.~1246--1247) but not any country's calendar of holy days
or its dispensations; and the law of annual confession beyond what the book cites.
\item Present discipline on the use of the 1962 books.
\item The origins of the Advent wreath and similar domestic customs, and the history of votive
Masses of Our Lady in Advent, for which no adequate sources were in hand.
\end{itemize}

\subsection{Method}

\textit{Sources.} Every statement of fact rests on an edition named in the References, read in
its printed text, in page images of it, or in a named electronic transcription. The stations and
readings of the 1962 Missal were read in a facsimile of the typical edition. The Latin of the two
collects quoted in full is Triptych's editorial text of the 1962 proper (the \textit{Roman 1962
Latin Proper Editorial Projection} of 3 September 2026), which gives the wording of the Pustet
Missal of 1862 read against the 1962 typical edition. The 1960 Code of Rubrics is cited from
\latin{Acta Apostolicae Sedis} 52. Its pages were read on page images, but whether each was a
facsimile of the printed volume or a later resetting of it has not been settled; the rubrics are
therefore cited by number and quoted only briefly, and a reader who needs their exact words should
consult the printed volume.

\textit{The Office of 1962.} No printed Breviary of 1961 or 1962 was read. The one-nocturn form of
Sunday Matins, with two lessons of Scripture and the opening of the homily (nn.~162\,a, 220), and
the three homily lessons of the Ember days (n.~222\,a), are stated from the rubrics. The English of
the homilies is Bute's, from his translation of the Breviary of his day, which printed the older
nine-lesson Office; several of the lessons quoted from it, such as the sermons of the second
nocturn, are not read in the 1962 Breviary. The course of Isaiah lessons in the daily table is that
of the Breviary of 1888, which assigns the same chapters to the same days as the Breviary of 1942.
The daily table applies the rules for the first two lessons on feasts of the third class (nn.~169,
221\,a) and for the omission of an impeded day's Scripture (n.~213); it should still be checked
against a 1962 Breviary where one is at hand.

\textit{The Fathers.} The Latin of St Gregory's homilies on the Gospels is not quoted; his words
are given in Bute's English and in the English \textit{Catena Aurea}. Bute numbers one homily
differently; the book uses the standard numbering throughout. The \textit{Catena Aurea} sometimes
joins or condenses passages, and every quotation taken from it is marked ``as the Catena gives
it.'' Two passages, Gregory on the reed shaken by the wind and Augustine on the desire of the heart
in his exposition of the thirty-seventh Psalm, are paraphrased and not quoted. Where an English
quotation omits words within it, the omission is marked.

\textit{Translations.} Every English quotation is taken from a published translation, named where
it is quoted and in the References. The translations are older ones, chosen because their printed
texts could be read, checked and quoted exactly: Bute's Breviary for the lessons and collects, the Mount
Melleray translation for St Bernard, the Nicene and Post-Nicene Fathers for Cyril, Augustine and
Leo, the Oxford translation of the \textit{Catena Aurea}, Hefele's English for the councils, Dalton for Gregory of Tours, Waterworth for Trent, Donovan
for the Roman Catechism, Mullan for St Ignatius, the Pustet English for St Francis de Sales, the
Raccolta of 1910 for the Angelus, Shepherd's English of Guéranger, and the Douay-Rheims Bible in
Challoner's revision for verses of Scripture. No new translation of any
text has been made for this book. Where the exposition explains a Latin word or phrase, the
explanation is marked ``gloss'' and is not offered as a text for prayer.

\textit{Spelling.} Quotations keep the spelling, punctuation and capitals of their editions, except
that the large capital letters with which some editions begin a passage are set in ordinary type,
and the signs for versicle and response are printed V. and R. The Latin collects taken from
Triptych's editorial text are a further exception: that text joins the printed lines, removes the
stress accents, prints the ae and oe ligatures as separate letters, writes consonantal i as j, and
takes its punctuation and the length of its conclusion from the 1962 edition.

\subsection{Law and its currency}

The book separates law, custom and counsel. Law is stated with its body, territory and date. The
1962 books belong to the time of the Code of Canon Law of 1917, in force from 19 May 1918 to 26
November 1983, and the book says what that Code required in Advent; it does not present those
requirements as binding now. Present law is stated from the Code of 1983, Paul VI's
\latin{Paenitemini} of 1966, and, for the United States, the bishops' pastoral statement of 18
November 1966 and their complementary norm of 1983 as amended in 1984, all as they stood on 29
November 2026. The English of the 1983 Code and of \latin{Paenitemini} is a working text; the Latin
governs. The American list of fasting days of 1889 is given as the particular law of that date; the
book does not establish whether its rule for the Fridays of Advent survived the Code of 1917. The
English Wednesday-and-Friday fast is known here only from a translator's note of 1870 and is
reported as a custom.

Customs are attributed to the time and place in which they are found. Indulgences recorded in the
Raccolta of 1910 are reported as the grants of that time and are not stated as current.

Counsel is the book's own. It is offered for prudent use, according to each reader's health, duties
and state of life, and it is best taken up with the advice of a confessor or spiritual director.

\subsection{Qualifications}

This book is not a liturgical book, an Ordo, a directory of law, or a substitute for spiritual
direction. It does not settle any question of rubrics or law for a particular church, and where
the law of a reader's diocese or country differs from what is stated here, that law governs.
